% GENERATED FILE. Edit canonical lessons, then run npm run generate:books.
% canonical-source-hash: fnv1a64:2cd30fee83988475
% canonical-lessons: JA-C179-koucha, JA-C179-osake, JA-C179-sosogu, JA-C179-kobosu, JA-C179-atatameru

\chapter{Tea, Sake and Spills}
\label{ch:ja-tea-sake-and-spills}
\hlchaptermodality{179}

\begin{chapteropening}
\textbf{By the end of this chapter:} \emph{I can talk about drinks and how I pour and warm them.}

Complete the last lesson of chapter 179: I can talk about drinks and how I pour and warm them.
\end{chapteropening}

\section[\texorpdfstring{\ja{こうちゃ} (kōcha)}{こうちゃ (kōcha)}]{\ja{こうちゃ} (kōcha) — black tea}

\label{lesson:JA-C179-koucha}



\begin{quote}
Before the new one: say the Japanese for milk, then the Japanese for to chill.
\end{quote}

\subsection*{You'll want to know: \ja{こうちゃ}}
\textbf{\ja{こうちゃ}} — \emph{kōcha} — \textquotedblleft{}black tea\textquotedblright{}.

\begin{quote}
\textbf{\ja{あさは} \ja{こうちゃを} \ja{のみます。}} — \emph{asa wa kōcha o nomimasu} — \textquotedblleft{}In the morning I drink black tea.\textquotedblright{}
\end{quote}

\begin{cousinweb}
\textbf{\ja{こうちゃ}} (\emph{kōcha}) is Sino-Japanese (kango): two Chinese characters read with their on'yomi, readings borrowed from Chinese. They are \emph{kō}, \textquotedblleft{}crimson\textquotedblright{} (Middle Chinese \emph{huwng}) and \emph{cha}, \textquotedblleft{}tea\textquotedblright{} (\emph{drae}) (KANJIDIC2; Middle Chinese in Baxter's transcription).
\end{cousinweb}

\begin{grammarlens}[title={What you've built}]
One more word for everyday life at A2.
\end{grammarlens}

\subsection*{Your turn}
\begin{itemize}
\raggedright
  \item \emph{Say it:} \emph{kōcha}
  \item \emph{Say it:} \emph{kōcha}, once more
  \item \emph{Say it:} \emph{asa wa kōcha o nomimasu} — In the morning I drink black tea.
  \item \emph{Recall:} say the Japanese for milk, then the Japanese for to chill, then say \emph{kōcha} again
\end{itemize}

\begin{tcolorbox}[breakable,colback=teal!4,colframe=teal!35!black,title={Before you move on}]
What does \ja{こうちゃ} mean? (\textquotedblleft{}Black tea\textquotedblright{}.) Say it once more.
\end{tcolorbox}
\section[\texorpdfstring{\ja{おさけ} (osake)}{おさけ (osake)}]{\ja{おさけ} (osake) — rice wine; sake, alcohol}

\label{lesson:JA-C179-osake}



\begin{quote}
Before the new one: say the Japanese for to chill, then the Japanese for black tea.
\end{quote}

\subsection*{You'll want to know: \ja{おさけ}}
\textbf{\ja{おさけ}} — \emph{osake} — \textquotedblleft{}rice wine; sake, alcohol\textquotedblright{}.

\begin{quote}
\textbf{\ja{ちちは} \ja{おさけを} \ja{すこし} \ja{のみます。}} — \emph{chichi wa osake o sukoshi nomimasu} — \textquotedblleft{}My father drinks a little sake.\textquotedblright{}
\end{quote}

\begin{cousinweb}
\textbf{\ja{おさけ}} (\emph{osake}) is built from polite \emph{o-} and native \emph{sake}, \textquotedblleft{}rice wine, alcohol\textquotedblright{}, inherited rather than borrowed; its kanji was attached later as a kun'yomi (JMdict, KANJIDIC2).
\end{cousinweb}

\begin{grammarlens}[title={What you've built}]
One more word for everyday life at A2.
\end{grammarlens}

\subsection*{Your turn}
\begin{itemize}
\raggedright
  \item \emph{Say it:} \emph{osake}
  \item \emph{Say it:} \emph{osake}, once more
  \item \emph{Say it:} \emph{chichi wa osake o sukoshi nomimasu} — My father drinks a little sake.
  \item \emph{Recall:} say the Japanese for to chill, then the Japanese for black tea, then say \emph{osake} again
\end{itemize}

\begin{tcolorbox}[breakable,colback=teal!4,colframe=teal!35!black,title={Before you move on}]
What does \ja{おさけ} mean? (\textquotedblleft{}Rice wine; sake, alcohol\textquotedblright{}.) Say it once more.
\end{tcolorbox}
\section[\texorpdfstring{\ja{そそぐ} (sosogu)}{そそぐ (sosogu)}]{\ja{そそぐ} (sosogu) — to pour; to sprinkle}

\label{lesson:JA-C179-sosogu}



\begin{quote}
Before the new one: say the Japanese for black tea, then the Japanese for rice wine.
\end{quote}

\subsection*{You'll want to know: \ja{そそぐ}}
\textbf{\ja{そそぐ}} — \emph{sosogu} — \textquotedblleft{}to pour; to sprinkle\textquotedblright{}.

\begin{quote}
\textbf{\ja{ゆのみに} \ja{おちゃを} \ja{そそぎます。}} — \emph{yunomi ni ocha o sosogimasu} — \textquotedblleft{}I pour tea into a teacup.\textquotedblright{}
\end{quote}

Polite \textbf{\ja{そそぎます}}, \emph{sosogimasu}; a godan verb.

\begin{cousinweb}
\textbf{\ja{そそぐ}} (\emph{sosogu}) is a native Japanese word (wago), inherited rather than borrowed from Chinese; its kanji was attached later as a kun'yomi (JMdict, KANJIDIC2).
\end{cousinweb}

\begin{grammarlens}[title={What you've built}]
One more word for everyday life at A2.
\end{grammarlens}

\subsection*{Your turn}
\begin{itemize}
\raggedright
  \item \emph{Say it:} \emph{sosogu}
  \item \emph{Say it:} \emph{sosogu}, once more
  \item \emph{Say it:} \emph{yunomi ni ocha o sosogimasu} — I pour tea into a teacup.
  \item \emph{Recall:} say the Japanese for black tea, then the Japanese for rice wine, then say \emph{sosogu} again
\end{itemize}

\begin{tcolorbox}[breakable,colback=teal!4,colframe=teal!35!black,title={Before you move on}]
What does \ja{そそぐ} mean? (\textquotedblleft{}To pour; to sprinkle\textquotedblright{}.) Say it once more.
\end{tcolorbox}
\section[\texorpdfstring{\ja{こぼす} (kobosu)}{こぼす (kobosu)}]{\ja{こぼす} (kobosu) — to spill; to drop}

\label{lesson:JA-C179-kobosu}



\begin{quote}
Before the new one: say the Japanese for rice wine, then the Japanese for to pour.
\end{quote}

\subsection*{You'll want to know: \ja{こぼす}}
\textbf{\ja{こぼす}} — \emph{kobosu} — \textquotedblleft{}to spill; to drop\textquotedblright{}.

\begin{quote}
\textbf{\ja{こどもが} \ja{ぎゅうにゅうを} \ja{こぼします。}} — \emph{kodomo ga gyūnyū o koboshimasu} — \textquotedblleft{}The child spills milk.\textquotedblright{}
\end{quote}

Polite \textbf{\ja{こぼします}}, \emph{koboshimasu}; a godan verb.

\begin{cousinweb}
\textbf{\ja{こぼす}} (\emph{kobosu}) is a native Japanese word (wago), inherited rather than borrowed from Chinese; its kanji was attached later as a kun'yomi (JMdict, KANJIDIC2).
\end{cousinweb}

\begin{grammarlens}[title={What you've built}]
One more word for everyday life at A2.
\end{grammarlens}

\subsection*{Your turn}
\begin{itemize}
\raggedright
  \item \emph{Say it:} \emph{kobosu}
  \item \emph{Say it:} \emph{kobosu}, once more
  \item \emph{Say it:} \emph{kodomo ga gyūnyū o koboshimasu} — The child spills milk.
  \item \emph{Recall:} say the Japanese for rice wine, then the Japanese for to pour, then say \emph{kobosu} again
\end{itemize}

\begin{tcolorbox}[breakable,colback=teal!4,colframe=teal!35!black,title={Before you move on}]
What does \ja{こぼす} mean? (\textquotedblleft{}To spill; to drop\textquotedblright{}.) Say it once more.
\end{tcolorbox}
\section[\texorpdfstring{\ja{あたためる} (atatameru)}{あたためる (atatameru)}]{\ja{あたためる} (atatameru) — to warm up; to heat}

\label{lesson:JA-C179-atatameru}



\begin{quote}
Before the new one: say the Japanese for to pour, then the Japanese for to spill.
\end{quote}

\subsection*{You'll want to know: \ja{あたためる}}
\textbf{\ja{あたためる}} — \emph{atatameru} — \textquotedblleft{}to warm up; to heat\textquotedblright{}.

\begin{quote}
\textbf{\ja{ぎゅうにゅうを} \ja{あたためます。}} — \emph{gyūnyū o atatamemasu} — \textquotedblleft{}I warm up the milk.\textquotedblright{}
\end{quote}

Polite \textbf{\ja{あたためます}}, \emph{atatamemasu}; an ichidan verb.

\begin{cousinweb}
\textbf{\ja{あたためる}} (\emph{atatameru}) is a native Japanese word (wago), inherited rather than borrowed from Chinese; its kanji was attached later as a kun'yomi, and it is built on the same stem as \emph{atatakai}, \textquotedblleft{}warm\textquotedblright{} (JMdict, KANJIDIC2).
\end{cousinweb}

\begin{grammarlens}[title={What you've built}]
One more word for everyday life at A2.
\end{grammarlens}

\subsection*{Your turn}
\begin{itemize}
\raggedright
  \item \emph{Say it:} \emph{atatameru}
  \item \emph{Say it:} \emph{atatameru}, once more
  \item \emph{Say it:} \emph{gyūnyū o atatamemasu} — I warm up the milk.
  \item \emph{Recall:} say the Japanese for to pour, then the Japanese for to spill, then say \emph{atatameru} again
\end{itemize}

\begin{tcolorbox}[breakable,colback=teal!4,colframe=teal!35!black,title={Before you move on}]
What does \ja{あたためる} mean? (\textquotedblleft{}To warm up; to heat\textquotedblright{}.) Say it once more.
\end{tcolorbox}
